% Compile-ready section. Requires: enumitem.

\section{GlyphChain and Legal Application-Account Recovery}
\label{sec:glyphchain-legal-application-account-recovery}

\noindent\textbf{Status:} The local protocol foundation, standalone monotonic
registry service and signed AION desktop-session enforcement are implemented.
An independently operated registry deployment, connection of the remaining
session issuers and live clean-device recovery and revocation exercises remain
production release gates.

\subsection{Purpose and Separation of Authority}

Application-account recovery is intentionally separate from whole-device,
Legal-key and GlyphChain network-identity backups.  Restoring files from a lost
machine does not prove that the person performing the restoration owns the
Tessaris account and does not grant validator authority.

The implemented protocol creates two distinct artefacts:

\begin{enumerate}[leftmargin=*]
  \item an encrypted account-root recovery kit, which must be retained outside
        the Tessaris runtime; and
  \item a root-signed public account descriptor, which may be published or
        retained independently without disclosing account secrets.
\end{enumerate}

The stable account identifier is derived from the Ed25519 account-root public
key.  The kit contains the corresponding private root encrypted with scrypt
using \(N=32768\), \(r=8\) and \(p=1\), followed by AES--256--GCM authenticated
encryption.  The recovery password is never written into the kit or runtime
state.  Production root-key custody uses the platform Keychain; development
tests use an owner-only local file.

The account-root key is not a Legal proof-device key, peer-to-peer transport
key or validator-finality key.  Deployment-managed validator-finality
credentials remain in approved secret custody and outside the recovery kit.
Account recovery neither restores nor copies those credentials and does not
grant validator participation rights.

\subsection{Clean-Device Recovery}

Recovery requires all of the following:

\begin{itemize}[leftmargin=*]
  \item the off-device encrypted recovery kit;
  \item the kit password;
  \item the expected account identifier;
  \item the current independently retained public descriptor;
  \item a new device identifier; and
  \item the exact destructive-recovery confirmation phrase.
\end{itemize}

Tessaris authenticates the kit, derives and verifies the root public key,
checks the expected account and workspace, verifies the current descriptor's
root signature, advances the monotonic recovery epoch and signs a replacement
descriptor.  The replacement lists only the newly recovered device and treats
every unlisted device as revoked.  Its continuity field commits to the exact
digest of the preceding descriptor.

Any consumer of this authority can reject a stale recovery epoch and a device
omitted from the active descriptor.  Wrong passwords, modified kits, modified
descriptors, unexpected account identifiers, recovery-kit destinations inside
the runtime and recovery into an already initialised target all fail closed.

\subsection{Authenticated Session Binding}

When an account descriptor is present, the installed Tessaris desktop signer
and backend verifier use the domain-separated
\texttt{aion.flow.desktop-session.v2} format.  Each signed request binds:

\begin{itemize}[leftmargin=*]
  \item the account identifier and recovery epoch;
  \item the authorised device, person and workspace;
  \item the exact HTTP method and target;
  \item the request-body digest; and
  \item a timestamp and single-use nonce.
\end{itemize}

The verifier rejects a pre-recovery epoch, a revoked or unknown device, a
mismatched account or workspace and replay of the nonce.  Account headers sent
to a backend without the corresponding descriptor fail closed instead of
falling back to the legacy session format.  Version~1 remains only as a
migration path when no account descriptor exists.  The remaining application
session issuers must adopt the same account, epoch and device binding before
production account recovery can be claimed.

The interactive operator command
\texttt{backend/scripts/application\_account\_recovery.py} exposes
\texttt{initialize}, \texttt{recover} and \texttt{status} operations.  It reads
passwords interactively and never accepts them as command-line arguments.  It
can write the public descriptor for independent publication while leaving the
encrypted recovery kit under owner control.

\subsection{Freshness and Publication Boundary}

The recovery kit cannot appoint its own latest account state.  A holder of an
old valid descriptor can prove that the descriptor was valid at one time, but
a local verifier cannot determine by itself whether a later recovery has
superseded it.  Production therefore requires a separately operated
current-descriptor registry or another independently controlled monotonic
publication channel.

The implemented standalone registry authenticates the root signature, accepts
epoch one as the initial state and thereafter requires an exact one-epoch
advance tied to the current descriptor digest.  It preserves append-only
history and rejects rollback, same-epoch equivocation, epoch gaps, broken
continuity, workspace changes, tampering and configured capacity exhaustion.
The service stores no recovery secrets, bounds request bodies and refuses
non-loopback operation without TLS and an out-of-band publication token of at
least 32 characters.  The approved secret store supplies that token through
\texttt{TESSARIS\_ACCOUNT\_REGISTRY\_PUBLISH\_TOKEN}; it is never accepted as a
command-line argument and is not required to read a public descriptor.

The backend can use the separately running service through
\texttt{TESSARIS\_ACCOUNT\_DESCRIPTOR\_REGISTRY\_URL}.  A production URL must
use HTTPS.  The client stores a local anti-rollback memory but never treats it
as fresh authority when the live registry is unavailable.  Therefore a network
failure fails closed rather than silently authorising a cached device.

A service running beside the desktop is not operationally independent.  A live
clean-device exercise must prove that the replacement device becomes active
and that every lost or replaced device is rejected using the separately
operated registry's current descriptor.  Until that deployment, connection of
the remaining session issuers and the live exercise pass, Tessaris must
describe this capability as an implemented local recovery and registry
foundation rather than production application-account recovery.

\subsection{Verified Local Behaviour}

The focused regression verifies:

\begin{itemize}[leftmargin=*]
  \item encrypted off-device kit creation without plaintext private-key
        disclosure;
  \item root-signed public descriptors and stable account-identifier
        verification;
  \item rejection of wrong passwords and tampered kits or descriptors;
  \item expected-account and workspace binding;
  \item monotonic recovery-epoch advancement;
  \item rejection of an old device and stale session epoch;
  \item clean-target enforcement;
  \item exact descriptor-chain continuity after recovery;
  \item version~2 signed desktop-session binding and fail-closed backend
        validation;
  \item append-only monotonic descriptor publication with idempotent replay;
  \item rollback, equivocation, gap, tampering and unsafe-path rejection;
  \item bounded standalone registry API behavior and HTTPS-only production
        client configuration; and
  \item persistent rollback detection without stale-cache authorization.
\end{itemize}

These tests are local protocol evidence.  They are not evidence of independent
publication, live multi-device revocation, production availability or recovery
of validator-finality authority.
